\section{No-wrap constant-mode decomposition}
\label{sec:constant-mode}

The determinant support of the physical coefficient lies in
\[
 [-B_{\mathscr D}Q,B_{\mathscr D}Q]\cap\Z
\]
for a fixed \(B_{\mathscr D}\) \citep{WangTPC32}.  We retain the
strong physical padding used there and choose an integer \(q\) so that
\begin{equation}
 q>2B_{\mathscr D}Q+2,
 \qquad q\asymp Q.
 \label{eq:constant-no-wrap-q}
\end{equation}
Thus reduction modulo \(q\) is injective on the support and no
nonzero determinant bin is congruent to zero modulo \(q\).  Define
the padded cyclic coefficient
\begin{equation}
 \widetilde A_C(a)
 =\sum_{\substack{n\in\Z\\n\equiv a\;(\mathrm{mod}\ q)}}A_C(n),
 \qquad a\in\Z/q\Z.
 \label{eq:constant-padded-A}
\end{equation}
The no-wrap condition gives
\begin{equation}
 \sum_{a\; (q)}|\widetilde A_C(a)|^2=D_X,
 \qquad
 \sum_{a\; (q)}\widetilde A_C(a)=Z_X,
 \qquad
 \widetilde A_C(0)=A_C(0)=0.
 \label{eq:constant-padded-identities}
\end{equation}

\subsection{The orthogonal projection}

Equip \(\C^{\Z/q\Z}\) with the unnormalized inner product
\[
 \langle f,g\rangle_q
 =\sum_{a\; (q)}f(a)\overline{g(a)}.
\]
Let \(\one_q\) be the constant-one vector and define
\begin{equation}
 P_0f=\frac{\langle f,\one_q\rangle_q}{q}\one_q,
 \qquad
 P_\perp=I-P_0.
 \label{eq:constant-projections}
\end{equation}

\begin{theorem}[Exact constant-mode disintegration and coherent
reassembly]
\label{thm:constant-mode}
Set
\begin{equation}
 A_C^\circ=P_\perp\widetilde A_C
 =\widetilde A_C-\frac{Z_X}{q}\one_q.
 \label{eq:constant-centered-A}
\end{equation}
Then
\begin{align}
 \sum_{a\; (q)}A_C^\circ(a)&=0,
 \label{eq:constant-centered-sum}\\
 \widetilde A_C
 &=A_C^\circ+\frac{Z_X}{q}\one_q,
 \label{eq:constant-decomposition}\\
 \boxed{
 D_X
 =\|A_C^\circ\|_{\ell^2(\Z/q\Z)}^2
  +\frac{|Z_X|^2}{q}.}
 \label{eq:constant-energy-decomposition}
\end{align}
For the unnormalized finite Fourier transform
\[
 \widehat f(r)=\sum_{a\; (q)}f(a)e_q(-ra),
\]
one has
\begin{equation}
 \widehat{A_C^\circ}(0)=0,
 \qquad
 \widehat{A_C^\circ}(r)
 =\widehat A_{C,q}(r)\quad(r\ne0),
 \label{eq:constant-fourier-decomposition}
\end{equation}
and hence
\begin{equation}
 \boxed{
 \sum_{\substack{r\; (q)\\r\ne0}}
 |\widehat A_{C,q}(r)|^2
 =qD_X-|Z_X|^2
 =q\|A_C^\circ\|_2^2.}
 \label{eq:constant-nonzero-frequency-energy}
\end{equation}
The deleted determinant bin also gives the exact coherent identity
\begin{equation}
 \boxed{
 Z_X=-\sum_{\substack{r\; (q)\\r\ne0}}
       \widehat A_{C,q}(r).}
 \label{eq:constant-coherent-reassembly}
\end{equation}
\end{theorem}

\begin{proof}
By \eqref{eq:constant-padded-identities},
\[
 P_0\widetilde A_C=\frac{Z_X}{q}\one_q.
\]
The ranges of \(P_0\) and \(P_\perp\) are orthogonal, while
\(\|\one_q\|_2^2=q\).  Pythagoras therefore gives
\eqref{eq:constant-energy-decomposition}.  The Fourier transform of
\((Z_X/q)\one_q\) equals \(Z_X\) at \(r=0\) and zero otherwise.
This proves \eqref{eq:constant-fourier-decomposition}.  Finite
Parseval then gives
\eqref{eq:constant-nonzero-frequency-energy}.  Finally, Fourier
inversion at the cyclic coordinate \(a=0\), together with
\(\widetilde A_C(0)=0\), gives
\[
 0=\frac1q\left(
 Z_X+\sum_{\substack{r\; (q)\\r\ne0}}
 \widehat A_{C,q}(r)\right),
\]
which is \eqref{eq:constant-coherent-reassembly}.
\end{proof}

\begin{corollary}[What centering alone does and does not prove]
\label{cor:centering-no-cancellation}
Let
\[
 S_X=\#\supp A_C.
\]
For the literal coefficient,
\begin{equation}
 \boxed{
 |Z_X|^2\le S_XD_X\le(q-1)D_X\le qD_X.}
 \label{eq:constant-cauchy}
\end{equation}
The final \(qD_X\) bound is the projection bound for an unrestricted
cyclic vector, but it is not sharp for the physically padded
coefficient because the zero coordinate is absent.  The centering
identity alone produces no fixed-power improvement over
\eqref{eq:constant-cauchy}.  On the other hand,
\eqref{eq:constant-coherent-reassembly} shows that a genuinely
coherent estimate for all nonzero Fourier modes can estimate \(Z_X\);
merely declaring the centered zero Fourier mode to vanish cannot.
\end{corollary}

\begin{proof}
Cauchy--Schwarz on the exact support gives
\[
 |Z_X|^2
 =\left|\sum_{a\in\supp\widetilde A_C}
            \widetilde A_C(a)\right|^2
 \le S_XD_X.
\]
The strong padding and \(\widetilde A_C(0)=0\) give
\(S_X\le q-1\).  The last assertions follow from
\cref{eq:constant-energy-decomposition,eq:constant-coherent-reassembly}.
\end{proof}

\begin{remark}[Dependence on the ambient group]
The decomposition is canonical after the physical no-wrap modulus
\(q\) has
been chosen.  Centering only on the set
\(\supp A_C\) is not an invariant replacement: exact cancellations
can change that support, and adding zero bins changes its average.
Different admissible no-wrap values of \(q\) give different constant
vectors, although the physical sum \(Z_X\) itself is independent of
\(q\).
\end{remark}

\begin{remark}[Weaker no-wrap padding]
If one assumes only
\(q>\diam(\supp A_C)\), reduction remains injective but a nonzero
integer bin can lie in residue class zero.  In that case Fourier
inversion gives the more general identity
\[
 Z_X=q\widetilde A_C(0)
 -\sum_{\substack{r\; (q)\\r\ne0}}\widehat A_{C,q}(r).
\]
The strong physical choice \eqref{eq:constant-no-wrap-q} is what
identifies \(\widetilde A_C(0)\) with the deleted literal bin
\(A_C(0)=0\).
\end{remark}

\begin{example}[Deleting the bin \(n=0\) is not centering]
Fix any \(z\in\C\) and take an abstract coefficient supported at
\(n=1\) with \(A(1)=z\).  Then \(A(0)=0\), but
\(\widehat A(0)=z\) and \(D=|z|^2\).  This finite example is not a
model of the physical arithmetic; it isolates the logical
nonimplication between the two zero statements in
\cref{prop:two-zeros}.
\end{example}
